% This is used for pdf compilation from qmd.
% Any special LaTeX commands should be updated in parallel with mymacros.sty which is used for html compilation.
% All rows *might* need an HTML equivalent. Those at the bottom of this file are identical.
% This is how the two files are included in the _quarto.yaml header of the qmd file
%format:
%  html:
%    defaults: [resources/latex/global-maths-definitions.yaml]
%  pdf:
%   include-in-header: ["include-in-header.tex"]

% the 'defaults' above secretly points to mymacros.sty where they really live.

\usepackage{gensymb}
\usepackage{amsmath}
\usepackage{xcolor}

\usepackage[skins,breakable]{tcolorbox}

% Code to reduce width of code blocks in PDF to encourage line wrapping
\usepackage{fvextra}
      \DefineVerbatimEnvironment{Highlighting}{Verbatim}{
        breaklines,
        breakanywhere,
        breaknonspaceingroup,
        commandchars=\\\{\},
        fontsize=\small
      }
% Also line wrap in verbatim blocks
\RecustomVerbatimEnvironment{verbatim}{Verbatim}{breaklines,fontsize=\footnotesize}

\usepackage{environ}

% Below here is an old way to defining new custom coloured boxes.
% We use nice extensions now.

% \definecolor{goals-col}{RGB}{201, 232, 222}
% 
% % For defining new callout boxes
% % HTML version is in CSS file
% % Need to turn on goals and goals-header as environments using the latex-environment extension for quarto
% % This is designed to work for the following fenced div syntax:
% % :::{.goals}
% % ::::{.goals-header}
% % TITLE
% % ::::
% % Actual box contents
% % :::
% \newtcolorbox{goals}[1][]{notitle, top=0pt, boxsep=0pt, coltitle=black!90!white, colback=white, colframe=goals-col,#1}
% \NewEnviron{goals-header}{\tcbsubtitle[top=1mm, bottom=1mm]{\BODY}}
% % It works by creating a box with no title. Then no spacing. Then a subtitle defined later.

% Increasing spacing between rows of tables
\def\arraystretch{2}

% Syntax for highlighting my KEY TERMS
% HTML version is included in css for .term.
% These are created with Quarto inline code
% like [CONTENTS]{.term} by the latex-environment package
\newcommand{\term}[1]{\textbf{#1}}

% Added to do derivatives, custom commands for HTML
% read by MathJax are inserted via diff.macro
% must keep naming in sync with esdiff package
\usepackage{esdiff}

%%%%%%%%%%%%%%%%%%%%%%%%%%%%%%%%%%%%%%%%%%%%%%%%%%%%%%%%%%%%
%%% BELOW HERE ARE COMMANDS SHARED WITH THE HTML OUTPUT %%%%
%%%%%%%%%%%%%%%%%%%%%%%%%%%%%%%%%%%%%%%%%%%%%%%%%%%%%%%%%%%%
%%%% KEEP THESE UPDATED ALSO IN THE MYMACROS.STY FILE   %%%%
%%%%%%%%%%%%%%%%%%%%%%%%%%%%%%%%%%%%%%%%%%%%%%%%%%%%%%%%%%%%

% Adding some operator names missing by default
\DeclareMathOperator{\cosec}{cosec}
\DeclareMathOperator{\powop}{pow} %\newcommand{\pow}{\operatorname{pow}}
\definecolor{frog}{rgb}{1,0,0}
\newcommand{\dx}{\,\text{d}x}
\newcommand{\du}{\,\text{d}u}
\newcommand{\dy}{\,\text{d}y}
\newcommand{\dt}{\,\text{d}t}

